\begin{table}[H]
 \centering
 \caption{Numerical accuracy of the slow-transition sensitivity}
 \label{tab:rewrite-research-share-accuracy}
 \small
 \begin{tabular}{rrrrr}
 \toprule
 $\sigma$ & Final horizon & Mesh points & Dynamic residual & Horizon change\\
 \midrule
 0.90 & 6,063.3 & 836 & $9.96\times10^{-10}$ & $5.01\times10^{-8}$\\
 1.00 & 5,579.8 & 888 & $1.06\times10^{-9}$ & $3.74\times10^{-8}$\\
 1.10 & 5,464.1 & 903 & $1.17\times10^{-9}$ & $4.00\times10^{-8}$\\
 1.50 & 3,648.3 & 1,100 & $1.21\times10^{-9}$ & $1.87\times10^{-8}$\\
 \bottomrule
 \end{tabular}
 \par\smallskip
 \begin{minipage}{0.95\textwidth}\footnotesize
 Dynamic residuals are the largest independently reconstructed errors
 across the full-horizon and dense first-decade checks. Horizon change
 compares the last two solutions in detrended logarithmic coordinates
 over years 0--500. Error bounds are rounded upward.
 The unit-elastic initial candidate uses the stricter tolerance
 $2\times10^{-7}$, followed by $10^{-8}$ and $10^{-9}$. Event-continuity
 log gaps are below $1.34\times10^{-15}$. The unit-elastic annual expenditure
 ratio is $0.000637427339$ after both horizon extensions; its log change
 from the initial candidate is below $3.44\times10^{-8}$. Across all four
 elasticities that change is below $6.47\times10^{-7}$. The remaining difference
 from the empirical target is a calibration discrepancy, not a numerical
 error or a relaxation of equilibrium conditions. The $\sigma=1.50$ case
 passes both Hamiltonian-support grids and the analytical continuation
 bound; the other three cases pass the global concavity check.
 \end{minipage}
\end{table}
